\documentclass{beamer}
\usepackage{amsmath}
\usepackage{amssymb}

% Choose a theme
\usetheme{Gotham}
\usecolortheme{default}
\usefonttheme{serif}

\title[Lecture 12: Finite Elasticity]{Lecture 12: Finite Elasticity}
\author{David Al-Attar}
\institute{Michaelmas term 2024}
\date{}

\begin{document}

% --- INTRODUCTION ---

\begin{frame}
    \titlepage
\end{frame}

\begin{frame}{Outline and Motivation}
    The main goal of this lecture is to derive the equations of motion for an elastic body from first principles.
    \pause
    \begin{block}{Applications}
        These equations are fundamental for modeling phenomena like seismic wave propagation in geophysics.
    \end{block}
    \pause
    \begin{alertblock}{Key Concept: Finite Elasticity}
        We will develop an exact theory of elasticity without assuming that deformations are small. This provides a deeper understanding of the underlying physics.
    \end{alertblock}
\end{frame}

% --- KINEMATICS ---

\begin{frame}{Kinematics: Describing the Motion}
    We describe motion by mapping particles from a \textbf{reference body} to their current position.
    
    \begin{itemize}
        \item At time $t$, the body occupies a region $M_t \subset \mathbb{R}^3$.
        \item We use a fixed reference body, $M$. Each point $\mathbf{x} \in M$ represents a unique particle.
        \item The \textbf{motion} is a map $\varphi$ that gives the position of particle $\mathbf{x}$ at time $t$:
        $$ \mathbf{y} = \varphi(\mathbf{x}, t) $$
        Here, $\mathbf{y} \in M_t$ is the current position. The motion $\varphi$ is the fundamental field in this theory.
    \end{itemize}
    \pause
    We assume the motion is smooth and one-to-one (no tearing or fractures).
\end{frame}

\begin{frame}{Key Quantities: Velocity & Deformation Gradient}
    From the motion $\varphi(\mathbf{x}, t)$, we can define several important quantities.

    \begin{block}{Velocity}
        The velocity of a particle $\mathbf{x}$ is the time derivative of its trajectory:
        $$ v_{i}(\mathbf{x},t) = \frac{\partial\varphi_{i}}{\partial t}(\mathbf{x},t) $$
        This is a \textit{material} velocity (for a fixed particle), not a spatial one.
    \end{block}
    \pause
    
    \begin{block}{Deformation Gradient}
        The deformation gradient, $\mathbf{F}$, quantifies local deformation:
        $$ F_{ij} = \frac{\partial\varphi_{i}}{\partial x_{j}} $$
        It describes how the motion differs between nearby points.
    \end{block}
\end{frame}

\begin{frame}{Key Quantities: The Jacobian}
    \begin{block}{Jacobian ($J$)}
        The determinant of $\mathbf{F}$ is the Jacobian. It relates volume elements.
        $$ J(\mathbf{x},t) = \det \mathbf{F}(\mathbf{x},t) > 0 $$
        $J>0$ ensures that the orientation of the material is preserved (e.g., a right-handed basis maps to a right-handed basis).
    \end{block}
    \pause
    \begin{exampleblock}{Transforming Integrals}
        The Jacobian is crucial for changing variables between the reference and current bodies:
        $$ \int_{M_{t}}f(\mathbf{y})d^{3}\mathbf{y} = \int_{M}f[\varphi(\mathbf{x},t)]J(\mathbf{x},t)d^{3}\mathbf{x} $$
        This proves very useful within the theory.
    \end{exampleblock}
\end{frame}

% --- CONSERVATION LAWS AND ENERGY ---

\begin{frame}{Conservation of Mass}
    The principle that mass is conserved for any sub-body leads to a key relationship between densities.
    
    \begin{itemize}
        \item Let $\rho(\mathbf{y}, t)$ be the density at a point $\mathbf{y}$ in the current configuration $M_t$.
        \item Let $\rho_0(\mathbf{x})$ be a time-independent density defined in the reference body $M$.
    \end{itemize}
    \pause
    
    \begin{alertblock}{Mass Conservation Equation}
        The relationship between the current (\textit{spatial}) density and the \textit{referential} density is:
        $$ \rho_0(\mathbf{x}) = J(\mathbf{x}, t) \rho[\varphi(\mathbf{x}, t), t] $$
        This is a statement of mass conservation expressed in the reference body.
    \end{alertblock}
\end{frame}

\begin{frame}{Energy of an Elastic Body}
    \begin{block}{Kinetic Energy ($\mathcal{T}$)}
        The total kinetic energy is the integral of the kinetic energy density over the reference body:
        $$ \mathcal{T}(t) = \frac{1}{2}\int_{M} \rho_0(\mathbf{x}) ||\mathbf{v}(\mathbf{x}, t)||^2 d^3\mathbf{x} $$
    \end{block}
    \pause
    
    \begin{block}{Potential Energy ($\mathcal{V}$)}
        For a \textbf{simple hyperelastic body}, potential energy is stored due to deformation and depends locally on the deformation gradient $\mathbf{F}$:
        $$ \mathcal{V}(t) = \int_{M} W[\mathbf{x}, \mathbf{F}(\mathbf{x}, t)] d^3\mathbf{x} $$
        The function $W$ is the \textbf{strain energy function}.
    \end{block}
\end{frame}

% --- SYMMETRIES ---

\begin{frame}{Symmetries of Potential Energy (1/2)}
    The \textbf{principle of material frame indifference} states that physical laws should not depend on the observer's frame of reference. This imposes constraints on the strain energy function $W$.
    \pause
    \begin{block}{Symmetry 1: Time-Translation Invariance}
        The material's properties do not change over time, so physics should be the same regardless of when you start the clock.
        \begin{itemize}
            \item This requires that the potential energy density $W$ has no explicit dependence on time, $t$.
        \end{itemize}
    \end{block}
\end{frame}

\begin{frame}{Symmetries of Potential Energy (2/2)}
    \begin{block}{Symmetry 2: Invariance Under Rigid Body Motions}
        The potential energy should not change if the body is simply translated or rotated without changing its shape or volume.
        \pause
        \begin{itemize}
            \item Invariance under translation implies $W$ cannot depend on the motion $\varphi$ itself, only on the deformation gradient $\mathbf{F}$.
            \item Invariance under rotation requires that for any rotation matrix $\mathbf{Q}$:
            $$ W(\mathbf{x}, \mathbf{F}) = W(\mathbf{x}, \mathbf{QF}) $$
        \end{itemize}
    \end{block}
\end{frame}

% --- EQUATIONS OF MOTION ---

\begin{frame}{The Principle of Stationary Action}
    The equations of motion can be derived from a variational principle.

    \begin{block}{Lagrangian Density ($\mathcal{L}$)}
        The Lagrangian density is the kinetic minus the potential energy density:
        $$ \mathcal{L} = \frac{1}{2} \rho_0(\mathbf{x}) ||\mathbf{v}||^2 - W(\mathbf{x}, \mathbf{F}) $$
    \end{block}
    \pause
    
    \begin{block}{Action ($\mathcal{S}$)}
        The action is the integral of the Lagrangian density over the reference body and a time interval $[0, T]$:
        $$ \mathcal{S}(\varphi) = \int_0^T \int_M \mathcal{L} \, d^3\mathbf{x} \, dt $$
    \end{block}
\end{frame}

\begin{frame}{Hamilton's Principle}
    \begin{alertblock}{Hamilton's Principle}
        The true motion $\varphi(\mathbf{x}, t)$ is a stationary point of the action, $\delta\mathcal{S} = 0$, for all smooth variations $\delta\varphi$ that vanish at the start and end times, $t=0$ and $t=T$.
    \end{alertblock}
    \begin{center}
    % You can add an image here if you like, e.g.:
    % \includegraphics[width=0.6\textwidth]{path_variation.png}
    \end{center}
    This powerful principle turns a dynamics problem into a calculus of variations problem.
\end{frame}


\begin{frame}{Deriving the Equations of Motion}
    Applying Hamilton's principle ($\delta\mathcal{S}=0$) yields the Euler-Lagrange equations.

    \begin{block}{Euler-Lagrange Equations (General Form)}
        For a general Lagrangian density $\mathcal{L}(\varphi, \mathbf{v}, \mathbf{F})$:
        $$ \frac{\partial\mathcal{L}}{\partial\varphi_{i}} - \frac{\partial}{\partial t}\left(\frac{\partial\mathcal{L}}{\partial v_{i}}\right) - \frac{\partial}{\partial x_{j}}\left(\frac{\partial\mathcal{L}}{\partial F_{ij}}\right) = 0 $$
    \end{block}
    \pause
    \begin{block}{First Piola-Kirchhoff Stress Tensor}
        We define the stress tensor as the derivative of the strain energy with respect to the deformation gradient:
        $$ T_{ij} = \frac{\partial W}{\partial F_{ij}} $$
        It relates forces in the current body to areas in the reference body.
    \end{block}
\end{frame}


\begin{frame}{The Final Equation of Motion}
    Substituting the specific Lagrangian for an elastic body into the Euler-Lagrange equations gives the final result.

    \begin{itemize}
        \item The time derivative term becomes the rate of change of momentum: $\frac{\partial}{\partial t}\left(\frac{\partial\mathcal{L}}{\partial v_{i}}\right) = \rho_0 \frac{\partial v_i}{\partial t}$.
        \item The spatial derivative term becomes the divergence of stress: $\frac{\partial}{\partial x_{j}}\left(\frac{\partial\mathcal{L}}{\partial F_{ij}}\right) = -\frac{\partial T_{ij}}{\partial x_j}$.
    \end{itemize}
    \pause
    \begin{alertblock}{Equation of Motion for an Elastic Body}
         Putting it all together, we get:
        $$ \rho_0 \frac{\partial v_i}{\partial t} - \frac{\partial T_{ij}}{\partial x_j} = 0 $$
    \end{alertblock}
\end{frame}


\begin{frame}{Boundary Conditions}
    Hamilton's principle also provides the \textbf{natural boundary conditions} that must hold on the surface of the body, $\partial M$.
    
    \begin{block}{Natural Boundary Condition}
        The variation $\delta\mathcal{S}$ is zero only if the following holds on the boundary:
        $$ T_{ij} \hat{n}_j = 0 $$
        where $\hat{\mathbf{n}}$ is the outward unit normal vector to the surface of the reference body, $\partial M$.
    \end{block}
    \pause
    
    \begin{itemize}
        \item The vector $\mathbf{t} = \mathbf{T}\hat{\mathbf{n}}$ is the \textbf{traction}, or surface force per unit reference area.
        \item This condition means the body has a \textbf{traction-free} surface.
    \end{itemize}
\end{frame}

% --- SUMMARY ---

\begin{frame}{Summary: What You Need to Know}
    \begin{itemize}
        \item \textbf{Kinematics:} Understand the concepts of the reference body, current configuration, motion ($\varphi$), deformation gradient ($\mathbf{F}$), and Jacobian ($J$).
        \pause
        \item \textbf{Mass Conservation:} Be able to derive and use the relation between the spatial density $\rho$ and the referential density $\rho_0$.
        \pause
        \item \textbf{Energy:} State the expressions for the kinetic and potential energy of a simple hyperelastic body and explain the underlying assumptions.
        \pause
        \item \textbf{Derivation:} Use Hamilton's principle to derive the equations of motion and associated boundary conditions for finite elasticity.
    \end{itemize}
\end{frame}

\end{document}
